\documentclass[a4paper,twoside]{article}

\usepackage{morewrites}

\usepackage[svgnames,dvipsnames,x11names]{xcolor}
\usepackage{graphicx}
\usepackage[english]{babel}
\usepackage[colorlinks=true, linkcolor=NavyBlue, urlcolor=NavyBlue, citecolor=NavyBlue]{hyperref}
\usepackage[a4paper]{geometry}
\usepackage{ts-typesetting}
\usepackage{ts-fonts}
\usepackage{xcolor}
\usepackage{epigraph}
\usepackage{marginnote}
\usepackage{multicol}
\usepackage{progressbar}
\usepackage{graphicx}
\usepackage{tcolorbox}
\usepackage{fvextra}
\usepackage{amsmath}
\usepackage{float}
\usepackage{seqsplit}
\renewcommand*{\marginfont}{
\footnotesize
\sloppy
\emergencystretch=1em
\hyphenpenalty=50
\exhyphenpenalty=50
}
\newlength{\tsmarginparavail}
\newlength{\tsmarginparalt}
\newlength{\tsmarginparbuf}
\setlength{\tsmarginparbuf}{6mm}
\makeatletter
\AtBeginDocument{
\setlength{\tsmarginparavail}{\dimexpr\paperwidth-\textwidth-1in-\oddsidemargin-\marginparsep-\tsmarginparbuf\relax}
\ifdim\tsmarginparavail<\z@\setlength{\tsmarginparavail}{\z@}\fi
\if@twoside
\setlength{\tsmarginparalt}{\dimexpr1in+\evensidemargin-\marginparsep-\tsmarginparbuf\relax}
\ifdim\tsmarginparalt<\z@\setlength{\tsmarginparalt}{\z@}\fi
\ifdim\tsmarginparalt<\tsmarginparavail
\setlength{\tsmarginparavail}{\tsmarginparalt}
\fi
\fi
\ifdim\tsmarginparavail<\marginparwidth
\setlength{\marginparwidth}{\tsmarginparavail}
\fi
}
\makeatother
\usepackage{ts-code}

\usepackage{lastpage}
\usepackage{refcount}
\makeatletter
\def\ps@plain{
\let\@oddhead\@empty
\let\@evenhead\@empty
\def\@oddfoot{
\hfil
\ifnum\getpagerefnumber{LastPage}>1\relax
\thepage
\fi
\hfil}
\let\@evenfoot\@oddfoot
}
\makeatother

\title{Limitations}
\author{}\date{}
\begin{document}
\maketitle
\section{Highlighting with \tscodeinline{ucharclasses}}
The \tscodeinline{ucharclasses} package allows you to define font transitions based on Unicode character classes. However, when using it in combination with the \tscodeinline{soul} package for highlighting, there can be compatibility issues. For example, the following code may not work as expected:
\begin{tscode}[lang=latex, engine=pygments]
\begin{Verbatim}[commandchars=\\\{\},breaklines, breakanywhere, commandchars=\\\{\}]
\PY{k}{\PYZbs{}documentclass}\PY{n+nb}{\PYZob{}}article\PY{n+nb}{\PYZcb{}}
\PY{k}{\PYZbs{}usepackage}\PY{n+nb}{\PYZob{}}soul\PY{n+nb}{\PYZcb{}}
\PY{k}{\PYZbs{}usepackage}\PY{n+nb}{\PYZob{}}fontspec\PY{n+nb}{\PYZcb{}}
\PY{k}{\PYZbs{}usepackage}\PY{n+na}{[Latin]}\PY{n+nb}{\PYZob{}}ucharclasses\PY{n+nb}{\PYZcb{}}

\PY{k}{\PYZbs{}setmainfont}\PY{n+nb}{\PYZob{}}Latin Modern Roman\PY{n+nb}{\PYZcb{}}
\PY{k}{\PYZbs{}setDefaultTransitions}\PY{n+nb}{\PYZob{}}\PY{k}{\PYZbs{}rmfamily}\PY{n+nb}{\PYZcb{}}\PY{n+nb}{\PYZob{}}\PY{k}{\PYZbs{}rmfamily}\PY{n+nb}{\PYZcb{}}
\PY{k}{\PYZbs{}enableTransitionRules}
\PY{k}{\PYZbs{}AtBeginDocument}\PY{n+nb}{\PYZob{}}\PY{k}{\PYZbs{}enableTransitionRules}\PY{n+nb}{\PYZcb{}}

\PY{k}{\PYZbs{}begin}\PY{n+nb}{\PYZob{}}document\PY{n+nb}{\PYZcb{}}
\PY{k}{\PYZbs{}hl}\PY{n+nb}{\PYZob{}}Foobar\PY{n+nb}{\PYZcb{}}
\PY{k}{\PYZbs{}end}\PY{n+nb}{\PYZob{}}document\PY{n+nb}{\PYZcb{}}
\end{Verbatim}
\end{tscode}
TeXSmith is meant to be compatible with \tslogo{XeLaTeX} and \tslogo{LuaLaTeX}, but due to the way \tscodeinline{ucharclasses} handles font transitions, it may interfere with the highlighting functionality provided by \tscodeinline{soul}. When TeXSmith detects \tslogo{XeLaTeX} it now disables \tscodeinline{soul} and falls back to a plain yellow \tscodeinline{\textbackslash{}hl\{...\}} (no box). With \tslogo{LuaLaTeX} it loads \tscodeinline{lua-\allowbreak{}ul} instead of \tscodeinline{soul}, keeping underlines without tripping over \tscodeinline{ucharclasses}. Output therefore differs slightly between \tslogo{XeLaTeX} and \tslogo{LuaLaTeX}, but both builds succeed.
\end{document}
